% =================================================
% Ученический рабочий лист
% =================================================

\worksheettitle
  {Точка, прямая, отрезок и луч}
  {Геометрические объекты и их обозначения}

\studentnameblock

\begin{checkpointbox}[Входная разминка]
  \begin{enumerate}[label=\textbf{\arabic*.},itemsep=0.45em]
    \item Отметьте ниже две точки и обозначьте их буквами \(A\) и \(B\).
      \workarea{13mm}
    \item Какой инструмент нужен, чтобы провести прямую через две точки?
      Ответ: \answerblank[45mm].
  \end{enumerate}
\end{checkpointbox}

\section{Точка}

\pointsexamplefigure

\begin{fillbox}[Дополните]
  Точку изображают маленькой \answerblank[28mm] и обозначают
  \answerblank[36mm] латинской буквой.

  На рисунке отмечены точки \answerblank[42mm].
\end{fillbox}

\begin{practicebox}[Задание 1]
  Отметьте три точки, не лежащие на одной прямой. Обозначьте их буквами
  \(M\), \(N\) и \(K\).
  \workarea{20mm}
\end{practicebox}

\section{Прямая}

\lineexamplefigure

\begin{fillbox}[Наблюдение]
  Прямая не имеет ни \answerblank[28mm], ни \answerblank[28mm].
  На рисунке показана только \answerblank[45mm] прямой.
\end{fillbox}

\begin{propertybox}[Через две точки]
  Через любые две различные точки можно провести прямую, и притом только
  одну.
\end{propertybox}

\begin{notebox}[Как рисуем и обозначаем]
  У прямой не ставят стрелки. Её можно назвать одной строчной буквой —
  прямая \(a\), или двумя точками на ней — прямая \(AB\).
\end{notebox}

\incidenceexamplefigure

\begin{practicebox}[Задание 2. По рисунку]
  \begin{enumerate}[label=\alph*),itemsep=0.35em]
    \item Какие точки лежат на прямой \(a\)? \answerblank[45mm]
    \item Какая точка не лежит на прямой \(a\)? \answerblank[25mm]
    \item Через какие две отмеченные точки проходит прямая \(a\)?
      \answerblank[35mm]
  \end{enumerate}
\end{practicebox}

\section{Отрезок}

\segmentexamplefigure

\begin{fillbox}[Дополните определение]
  Точки \(A\) и \(B\) вместе со всеми точками прямой, расположенными между
  ними, образуют \answerblank[30mm] \(AB\).

  Точки \(A\) и \(B\) называют \answerblank[34mm] отрезка.
\end{fillbox}

\begin{checkpointbox}[Важное различие]
  Прямая не имеет концов, а отрезок имеет \answerblank[18mm] конца.
  Поэтому длину можно измерить у \answerblank[34mm], но не у всей прямой.
\end{checkpointbox}

\threelinepointsfigure

\begin{practicebox}[Задание 3. Найдите все отрезки]
  На прямой отмечены точки \(A\), \(B\) и \(C\). Назовите все отрезки,
  концами которых служат отмеченные точки:
  \[
    \answerblank[24mm],\qquad
    \answerblank[24mm],\qquad
    \answerblank[24mm].
  \]
\end{practicebox}

\section{Луч}

\rayexamplefigure

\begin{fillbox}[Дополните]
  Луч имеет \answerblank[28mm], но не имеет \answerblank[28mm].
  Луч \(AB\) начинается в точке \answerblank[15mm] и проходит через точку
  \answerblank[15mm].
\end{fillbox}

\begin{notebox}[Школьное обозначение]
  Пишем \textbf{луч \(AB\)} — без стрелки над буквами. Первая буква
  обозначает начало луча.
\end{notebox}

\oppositeraysfigure

\begin{practicebox}[Задание 4]
  По рисунку назовите:
  \begin{enumerate}[label=\alph*),itemsep=0.35em]
    \item луч, идущий из точки \(A\) через точку \(B\):
      \answerblank[26mm];
    \item луч, идущий из точки \(A\) через точку \(C\):
      \answerblank[26mm];
    \item общее начало этих лучей: \answerblank[18mm].
  \end{enumerate}
\end{practicebox}

\section{Различаем объекты}

\identifyobjectsfigure

\begin{fillbox}[Заполните таблицу]
  \begin{center}
    \renewcommand{\arraystretch}{1.42}
    \begin{tabularx}{\textwidth}{
      >{\raggedright\arraybackslash}p{0.20\textwidth}
      >{\centering\arraybackslash}p{0.18\textwidth}
      >{\raggedright\arraybackslash}p{0.28\textwidth}
      >{\centering\arraybackslash}X}
      \toprule
      \rowcolor{TableHeader}
      Объект & Сколько концов? & Как продолжается? & Можно измерить длину? \\
      \midrule
      Прямая & \answerblank[18mm] & \answerblank[35mm] & \answerblank[18mm] \\
      Отрезок & \answerblank[18mm] & \answerblank[35mm] & \answerblank[18mm] \\
      Луч & \answerblank[18mm] & \answerblank[35mm] & \answerblank[18mm] \\
      \bottomrule
    \end{tabularx}
  \end{center}
\end{fillbox}

\section{Построение}

\begin{practicebox}[Задание 5. Работайте линейкой]
  На рисунке ниже:
  \begin{enumerate}[label=\arabic*),itemsep=0.35em]
    \item проведите прямую \(AB\);
    \item проведите отрезок \(CD\);
    \item проведите луч \(DA\).
  \end{enumerate}
  \constructionpointsfigure
\end{practicebox}

\section{Верно или неверно?}

\begin{checkpointbox}[Задание 6]
  Поставьте знак «\(+\)», если утверждение верно, и «\(-\)», если неверно.
  \begin{enumerate}[label=\textbf{\arabic*.},itemsep=0.42em]
    \item \choicebox\ Прямая \(AB\) имеет концы \(A\) и \(B\).
    \item \choicebox\ Через две различные точки проходит только одна прямая.
    \item \choicebox\ Отрезки \(AB\) и \(BA\) — один и тот же отрезок.
    \item \choicebox\ Луч \(AB\) начинается в точке \(B\).
    \item \choicebox\ Лучи \(AB\) и \(BA\) обычно различны.
  \end{enumerate}
\end{checkpointbox}

\begin{reflectionbox}[Итог]
  Одним-двумя предложениями объясните, чем прямая, отрезок и луч отличаются
  друг от друга.
  \writinglines{1}
\end{reflectionbox}
